\section{Geometry and uncertainty}
A two-dimensional trajectory can be rotated through an angle $\theta$.
\begin{equation}
  R(\theta) =
  \begin{pmatrix}
    \cos\theta & -\sin\theta \\
    \sin\theta & \cos\theta
  \end{pmatrix}
  \label{eq:rotation_matrix}
\end{equation}
The squared length of a vector is preserved by this rotation.
\begin{equation}
  \|R(\theta)\vect{x}\|_2^2 = \vect{x}^\top\vect{x}
  \label{eq:length_preservation}
\end{equation}
For a random two-dimensional position, we use a covariance matrix with $|\rho|<1$.
\begin{equation}
  \Sigma = \begin{pmatrix}
    \sigma_x^2 & \rho\sigma_x\sigma_y \\
    \rho\sigma_x\sigma_y & \sigma_y^2
  \end{pmatrix}
  \label{eq:position_covariance}
\end{equation}
A standard normal density provides a compact inline-to-display example.
\begin{equation}
  p(z) = \frac{1}{\sqrt{2\pi}}\exp\!\left(-\frac{z^2}{2}\right)
  \label{eq:normal_density}
\end{equation}
Its probability mass over an interval is the corresponding definite integral.
\begin{equation}
  \Pr(a\leq Z\leq b) = \int_a^b p(z)\,dz
  \label{eq:interval_probability}
\end{equation}
